% !TEX root = ../main.tex
% =====================================================================
% NOT USED ANY MORE -- kept on purpose, do not delete.
%
% This was a standalone figure for the bottleneck. It is now panel (a) of
% figures/fig_variants_tikz.tex, so that all seven encoder variants sit in
% one figure. Keep this file in case the standalone version is ever wanted
% back: % !TEX root = ../main.tex
% =====================================================================
% NOT USED ANY MORE -- kept on purpose, do not delete.
%
% This was a standalone figure for the bottleneck. It is now panel (a) of
% figures/fig_variants_tikz.tex, so that all seven encoder variants sit in
% one figure. Keep this file in case the standalone version is ever wanted
% back: % !TEX root = ../main.tex
% =====================================================================
% NOT USED ANY MORE -- kept on purpose, do not delete.
%
% This was a standalone figure for the bottleneck. It is now panel (a) of
% figures/fig_variants_tikz.tex, so that all seven encoder variants sit in
% one figure. Keep this file in case the standalone version is ever wanted
% back: % !TEX root = ../main.tex
% =====================================================================
% NOT USED ANY MORE -- kept on purpose, do not delete.
%
% This was a standalone figure for the bottleneck. It is now panel (a) of
% figures/fig_variants_tikz.tex, so that all seven encoder variants sit in
% one figure. Keep this file in case the standalone version is ever wanted
% back: \input{figures/fig_bottleneck_tikz} inside a figure environment is
% all it needs.
% =====================================================================
% fig_bottleneck_tikz.tex -- the squeeze/expand trick, drawn.
%
% Source of truth: seanet/features.py, class MultiScaleSepBottleneckBlock
%   "reduce":  nn.Conv1d(d, r, kernel_size=1)      d -> r
%   "expand":  nn.Conv1d(r, d, kernel_size=1)      r -> d
%   r = max(1, d // bottleneck_ratio), ratio = 4
%
% The point of the picture: input and output are IDENTICAL in both rows.
% Only the middle is narrow. The box heights are drawn to scale (24 mm for
% d, 6 mm for r = d/4), so "narrow middle" is something you can see and not
% just something you read.
% =====================================================================

\begin{tikzpicture}[
  font=\small,
  >={Stealth[length=2mm]},
  % a stack of channels: the height IS the channel count, drawn to scale
  chan/.style    = {draw=black!55, fill=black!10, rounded corners=1pt, minimum width=7mm},
  % a 1x1 convolution
  convbox/.style = {draw=orange!65!black, fill=orange!18, rounded corners=2pt,
                    minimum height=9mm, align=center, inner xsep=5pt},
  cost/.style    = {black!75, align=center},
]

% =====================================================================
% (a) the standard pointwise convolution: one big d -> d step
% =====================================================================
\node[chan, minimum height=24mm] (a1) at (0,0) {};
\node[convbox]                   (ac) at (5,0) {$1{\times}1$ pointwise\\[-1pt]$\dmodel \to \dmodel$};
\node[chan, minimum height=24mm] (a2) at (10,0) {};

\draw[->] (a1.east) -- (ac.west);
\draw[->] (ac.east) -- (a2.west);

\node[anchor=east, black!70] at (a1.west) {$\dmodel = 64$\ };
\node[anchor=west, black!70] at (a2.east) {\ $\dmodel = 64$};

\node[anchor=south, black!80] at (5,1.55) {\textbf{(a) the standard block}};
\node[cost, anchor=north] at (5,-1.5)
     {cost $= \dmodel \cdot \dmodel = \dmodel^2 = 4096$ weights};

% =====================================================================
% (b) the bottleneck: squeeze down to r, then expand back up
% =====================================================================
\begin{scope}[yshift=-4.4cm]

\node[chan, minimum height=24mm] (b1) at (0,0) {};
\node[convbox]                   (bs) at (2.7,0) {squeeze\\[-1pt]$\dmodel \to r$};
\node[chan, minimum height=6mm]  (bm) at (5,0) {};      % 6mm = 24mm / 4, drawn to scale
\node[convbox]                   (be) at (7.3,0) {expand\\[-1pt]$r \to \dmodel$};
\node[chan, minimum height=24mm] (b2) at (10,0) {};

\draw[->] (b1.east) -- (bs.west);
\draw[->] (bs.east) -- (bm.west);
\draw[->] (bm.east) -- (be.west);
\draw[->] (be.east) -- (b2.west);

\node[anchor=east, black!70] at (b1.west) {$\dmodel = 64$\ };
\node[anchor=west, black!70] at (b2.east) {\ $\dmodel = 64$};
\node[anchor=south, black!70] at (bm.north) {$r = \dmodel/4 = 16$};

\node[anchor=south, black!80] at (5,1.55)
     {\textbf{(b) the bottleneck block} -- what \seanet{} uses};
\node[cost, anchor=north] at (5,-1.5)
     {cost $= \underbrace{\dmodel \cdot r}_{\text{squeeze}} +
              \underbrace{r \cdot \dmodel}_{\text{expand}}
             = 2\dmodel r = \dfrac{\dmodel^2}{2} = 2048$ weights};

\end{scope}

% ---- the take-away, tying the two rows together ----------------------
\node[draw=black!45, fill=black!4, rounded corners=2pt, inner sep=5pt, align=center,
      anchor=north] at (5,-7.5)
     {Same input, same output -- only the middle is narrow.\\
      \textbf{Half the weights, for any $\dmodel$.}};

\end{tikzpicture}
 inside a figure environment is
% all it needs.
% =====================================================================
% fig_bottleneck_tikz.tex -- the squeeze/expand trick, drawn.
%
% Source of truth: seanet/features.py, class MultiScaleSepBottleneckBlock
%   "reduce":  nn.Conv1d(d, r, kernel_size=1)      d -> r
%   "expand":  nn.Conv1d(r, d, kernel_size=1)      r -> d
%   r = max(1, d // bottleneck_ratio), ratio = 4
%
% The point of the picture: input and output are IDENTICAL in both rows.
% Only the middle is narrow. The box heights are drawn to scale (24 mm for
% d, 6 mm for r = d/4), so "narrow middle" is something you can see and not
% just something you read.
% =====================================================================

\begin{tikzpicture}[
  font=\small,
  >={Stealth[length=2mm]},
  % a stack of channels: the height IS the channel count, drawn to scale
  chan/.style    = {draw=black!55, fill=black!10, rounded corners=1pt, minimum width=7mm},
  % a 1x1 convolution
  convbox/.style = {draw=orange!65!black, fill=orange!18, rounded corners=2pt,
                    minimum height=9mm, align=center, inner xsep=5pt},
  cost/.style    = {black!75, align=center},
]

% =====================================================================
% (a) the standard pointwise convolution: one big d -> d step
% =====================================================================
\node[chan, minimum height=24mm] (a1) at (0,0) {};
\node[convbox]                   (ac) at (5,0) {$1{\times}1$ pointwise\\[-1pt]$\dmodel \to \dmodel$};
\node[chan, minimum height=24mm] (a2) at (10,0) {};

\draw[->] (a1.east) -- (ac.west);
\draw[->] (ac.east) -- (a2.west);

\node[anchor=east, black!70] at (a1.west) {$\dmodel = 64$\ };
\node[anchor=west, black!70] at (a2.east) {\ $\dmodel = 64$};

\node[anchor=south, black!80] at (5,1.55) {\textbf{(a) the standard block}};
\node[cost, anchor=north] at (5,-1.5)
     {cost $= \dmodel \cdot \dmodel = \dmodel^2 = 4096$ weights};

% =====================================================================
% (b) the bottleneck: squeeze down to r, then expand back up
% =====================================================================
\begin{scope}[yshift=-4.4cm]

\node[chan, minimum height=24mm] (b1) at (0,0) {};
\node[convbox]                   (bs) at (2.7,0) {squeeze\\[-1pt]$\dmodel \to r$};
\node[chan, minimum height=6mm]  (bm) at (5,0) {};      % 6mm = 24mm / 4, drawn to scale
\node[convbox]                   (be) at (7.3,0) {expand\\[-1pt]$r \to \dmodel$};
\node[chan, minimum height=24mm] (b2) at (10,0) {};

\draw[->] (b1.east) -- (bs.west);
\draw[->] (bs.east) -- (bm.west);
\draw[->] (bm.east) -- (be.west);
\draw[->] (be.east) -- (b2.west);

\node[anchor=east, black!70] at (b1.west) {$\dmodel = 64$\ };
\node[anchor=west, black!70] at (b2.east) {\ $\dmodel = 64$};
\node[anchor=south, black!70] at (bm.north) {$r = \dmodel/4 = 16$};

\node[anchor=south, black!80] at (5,1.55)
     {\textbf{(b) the bottleneck block} -- what \seanet{} uses};
\node[cost, anchor=north] at (5,-1.5)
     {cost $= \underbrace{\dmodel \cdot r}_{\text{squeeze}} +
              \underbrace{r \cdot \dmodel}_{\text{expand}}
             = 2\dmodel r = \dfrac{\dmodel^2}{2} = 2048$ weights};

\end{scope}

% ---- the take-away, tying the two rows together ----------------------
\node[draw=black!45, fill=black!4, rounded corners=2pt, inner sep=5pt, align=center,
      anchor=north] at (5,-7.5)
     {Same input, same output -- only the middle is narrow.\\
      \textbf{Half the weights, for any $\dmodel$.}};

\end{tikzpicture}
 inside a figure environment is
% all it needs.
% =====================================================================
% fig_bottleneck_tikz.tex -- the squeeze/expand trick, drawn.
%
% Source of truth: seanet/features.py, class MultiScaleSepBottleneckBlock
%   "reduce":  nn.Conv1d(d, r, kernel_size=1)      d -> r
%   "expand":  nn.Conv1d(r, d, kernel_size=1)      r -> d
%   r = max(1, d // bottleneck_ratio), ratio = 4
%
% The point of the picture: input and output are IDENTICAL in both rows.
% Only the middle is narrow. The box heights are drawn to scale (24 mm for
% d, 6 mm for r = d/4), so "narrow middle" is something you can see and not
% just something you read.
% =====================================================================

\begin{tikzpicture}[
  font=\small,
  >={Stealth[length=2mm]},
  % a stack of channels: the height IS the channel count, drawn to scale
  chan/.style    = {draw=black!55, fill=black!10, rounded corners=1pt, minimum width=7mm},
  % a 1x1 convolution
  convbox/.style = {draw=orange!65!black, fill=orange!18, rounded corners=2pt,
                    minimum height=9mm, align=center, inner xsep=5pt},
  cost/.style    = {black!75, align=center},
]

% =====================================================================
% (a) the standard pointwise convolution: one big d -> d step
% =====================================================================
\node[chan, minimum height=24mm] (a1) at (0,0) {};
\node[convbox]                   (ac) at (5,0) {$1{\times}1$ pointwise\\[-1pt]$\dmodel \to \dmodel$};
\node[chan, minimum height=24mm] (a2) at (10,0) {};

\draw[->] (a1.east) -- (ac.west);
\draw[->] (ac.east) -- (a2.west);

\node[anchor=east, black!70] at (a1.west) {$\dmodel = 64$\ };
\node[anchor=west, black!70] at (a2.east) {\ $\dmodel = 64$};

\node[anchor=south, black!80] at (5,1.55) {\textbf{(a) the standard block}};
\node[cost, anchor=north] at (5,-1.5)
     {cost $= \dmodel \cdot \dmodel = \dmodel^2 = 4096$ weights};

% =====================================================================
% (b) the bottleneck: squeeze down to r, then expand back up
% =====================================================================
\begin{scope}[yshift=-4.4cm]

\node[chan, minimum height=24mm] (b1) at (0,0) {};
\node[convbox]                   (bs) at (2.7,0) {squeeze\\[-1pt]$\dmodel \to r$};
\node[chan, minimum height=6mm]  (bm) at (5,0) {};      % 6mm = 24mm / 4, drawn to scale
\node[convbox]                   (be) at (7.3,0) {expand\\[-1pt]$r \to \dmodel$};
\node[chan, minimum height=24mm] (b2) at (10,0) {};

\draw[->] (b1.east) -- (bs.west);
\draw[->] (bs.east) -- (bm.west);
\draw[->] (bm.east) -- (be.west);
\draw[->] (be.east) -- (b2.west);

\node[anchor=east, black!70] at (b1.west) {$\dmodel = 64$\ };
\node[anchor=west, black!70] at (b2.east) {\ $\dmodel = 64$};
\node[anchor=south, black!70] at (bm.north) {$r = \dmodel/4 = 16$};

\node[anchor=south, black!80] at (5,1.55)
     {\textbf{(b) the bottleneck block} -- what \seanet{} uses};
\node[cost, anchor=north] at (5,-1.5)
     {cost $= \underbrace{\dmodel \cdot r}_{\text{squeeze}} +
              \underbrace{r \cdot \dmodel}_{\text{expand}}
             = 2\dmodel r = \dfrac{\dmodel^2}{2} = 2048$ weights};

\end{scope}

% ---- the take-away, tying the two rows together ----------------------
\node[draw=black!45, fill=black!4, rounded corners=2pt, inner sep=5pt, align=center,
      anchor=north] at (5,-7.5)
     {Same input, same output -- only the middle is narrow.\\
      \textbf{Half the weights, for any $\dmodel$.}};

\end{tikzpicture}
 inside a figure environment is
% all it needs.
% =====================================================================
% fig_bottleneck_tikz.tex -- the squeeze/expand trick, drawn.
%
% Source of truth: seanet/features.py, class MultiScaleSepBottleneckBlock
%   "reduce":  nn.Conv1d(d, r, kernel_size=1)      d -> r
%   "expand":  nn.Conv1d(r, d, kernel_size=1)      r -> d
%   r = max(1, d // bottleneck_ratio), ratio = 4
%
% The point of the picture: input and output are IDENTICAL in both rows.
% Only the middle is narrow. The box heights are drawn to scale (24 mm for
% d, 6 mm for r = d/4), so "narrow middle" is something you can see and not
% just something you read.
% =====================================================================

\begin{tikzpicture}[
  font=\small,
  >={Stealth[length=2mm]},
  % a stack of channels: the height IS the channel count, drawn to scale
  chan/.style    = {draw=black!55, fill=black!10, rounded corners=1pt, minimum width=7mm},
  % a 1x1 convolution
  convbox/.style = {draw=orange!65!black, fill=orange!18, rounded corners=2pt,
                    minimum height=9mm, align=center, inner xsep=5pt},
  cost/.style    = {black!75, align=center},
]

% =====================================================================
% (a) the standard pointwise convolution: one big d -> d step
% =====================================================================
\node[chan, minimum height=24mm] (a1) at (0,0) {};
\node[convbox]                   (ac) at (5,0) {$1{\times}1$ pointwise\\[-1pt]$\dmodel \to \dmodel$};
\node[chan, minimum height=24mm] (a2) at (10,0) {};

\draw[->] (a1.east) -- (ac.west);
\draw[->] (ac.east) -- (a2.west);

\node[anchor=east, black!70] at (a1.west) {$\dmodel = 64$\ };
\node[anchor=west, black!70] at (a2.east) {\ $\dmodel = 64$};

\node[anchor=south, black!80] at (5,1.55) {\textbf{(a) the standard block}};
\node[cost, anchor=north] at (5,-1.5)
     {cost $= \dmodel \cdot \dmodel = \dmodel^2 = 4096$ weights};

% =====================================================================
% (b) the bottleneck: squeeze down to r, then expand back up
% =====================================================================
\begin{scope}[yshift=-4.4cm]

\node[chan, minimum height=24mm] (b1) at (0,0) {};
\node[convbox]                   (bs) at (2.7,0) {squeeze\\[-1pt]$\dmodel \to r$};
\node[chan, minimum height=6mm]  (bm) at (5,0) {};      % 6mm = 24mm / 4, drawn to scale
\node[convbox]                   (be) at (7.3,0) {expand\\[-1pt]$r \to \dmodel$};
\node[chan, minimum height=24mm] (b2) at (10,0) {};

\draw[->] (b1.east) -- (bs.west);
\draw[->] (bs.east) -- (bm.west);
\draw[->] (bm.east) -- (be.west);
\draw[->] (be.east) -- (b2.west);

\node[anchor=east, black!70] at (b1.west) {$\dmodel = 64$\ };
\node[anchor=west, black!70] at (b2.east) {\ $\dmodel = 64$};
\node[anchor=south, black!70] at (bm.north) {$r = \dmodel/4 = 16$};

\node[anchor=south, black!80] at (5,1.55)
     {\textbf{(b) the bottleneck block} -- what \seanet{} uses};
\node[cost, anchor=north] at (5,-1.5)
     {cost $= \underbrace{\dmodel \cdot r}_{\text{squeeze}} +
              \underbrace{r \cdot \dmodel}_{\text{expand}}
             = 2\dmodel r = \dfrac{\dmodel^2}{2} = 2048$ weights};

\end{scope}

% ---- the take-away, tying the two rows together ----------------------
\node[draw=black!45, fill=black!4, rounded corners=2pt, inner sep=5pt, align=center,
      anchor=north] at (5,-7.5)
     {Same input, same output -- only the middle is narrow.\\
      \textbf{Half the weights, for any $\dmodel$.}};

\end{tikzpicture}
